\chapter{Fundamentação teórica}
\label{cap:fundamentacao-teorica}

[Apresente os conceitos e trabalhos que fundamentam a pesquisa.]

\section{Citações}

% Citações no sistema autor-data (NBR 10520:2023):
%  \cite{chave}       -> (Sobrenome, ano)
%  \cite[p. 10]{chave} -> (Sobrenome, ano, p. 10)
%  \citeonline{chave} -> Sobrenome (ano)
%  \citeonline[p. 10]{chave} -> Sobrenome (ano, p. 10)

Citação indireta com o autor fora dos parênteses: segundo \citeonline{angeloni2008}, [texto parafraseado]. Citação indireta com o autor entre parênteses: [texto parafraseado] \cite{leite2013}. Várias obras na mesma chamada são separadas por ponto e vírgula \cite{angeloni2008,oliveira2013}.

Citação direta com até três linhas, entre aspas duplas: de acordo com \citeonline[p. 25]{angeloni2008}, ``[trecho transcrito do documento consultado]''.

Citação direta com mais de três linhas, com recuo de 4 cm, fonte menor, espaço simples e sem aspas:

\begin{citacao}
Exemplo de citação direta longa: trecho transcrito com mais de três linhas, apresentado com recuo de 4 cm da margem esquerda, em fonte menor, espaço simples e sem aspas. Substitua este parágrafo pelo trecho transcrito do documento consultado, mantendo a chamada da fonte ao final \cite[p. 160]{oliveira2013}.
\end{citacao}

Documentos de pessoa jurídica são citados pelo nome completo ou pela sigla da instituição, como \cite{abnt2023}.

\section{Equações}

% Equações destacadas no texto e numeradas com algarismos arábicos entre
% parênteses, alinhados à direita (NBR 14724).

Exemplo de equação numerada, referenciada no texto como Equação~\eqref{eq:exemplo}:

\begin{equation}
	y = a x^{2} + b x + c
	\label{eq:exemplo}
\end{equation}
em que $a$, $b$ e $c$ são constantes reais, com $a \neq 0$.
